\documentclass[12pt]{report}
\usepackage[english, russian]{babel}
\usepackage{url}
\usepackage[utf8x]{inputenc}
\usepackage{amsmath, amsfonts}
\usepackage{graphicx}
\usepackage{parskip}
\usepackage{fancyhdr}
\usepackage{vmargin}
\usepackage{upgreek}
\usepackage{multicol}
\usepackage{bm}
\usepackage{xcolor}
\usepackage[colorlinks=true, linkcolor=blue, urlcolor=blue, citecolor=blue]{hyperref}
\usepackage{enumitem}
\usepackage{fancyvrb}
\usepackage{setspace}

\fancyhead[LE]{\leftmark}
\fancyhead[RO]{}
\fancyfoot[C]{\thepage}

\usepackage{pgfplots}

\pgfplotsset{compat=1.18}

\makeatletter
\AddEnumerateCounter{\asbuk}{\russian@alph}{щ}
\makeatother

\setcounter{tocdepth}{3}
\setcounter{secnumdepth}{3}

\let\tau\uptau
\let\eps\varepsilon
\let\gam\gamma

\setmarginsrb{3 cm}{2.5 cm}{3 cm}{2.5 cm}{1 cm}{1.5 cm}{1 cm}{1.5 cm}

\pagestyle{fancy}

%%%%%%%%%%%%%%%%%%%
\begin{document}
\begin{titlepage}

\vfill
\begin{center}

\rule{\linewidth}{0.5 mm} \\[0.1 cm]
{\huge \bfseries Оптимальный период передачи\\[0.01cm]

дискретного сигнала между\\[0.01cm]

двумя системами\\[0.01cm]
}
\rule{\linewidth}{0.5 mm}
\\[2.0 cm]

{\Large Теоретические основы алгоритмов и их эмпирическая верификация.}\\[5.0 cm]

    \begin{flushleft}
    \large \emph{Автор: \\Тубольцев М. Е.}
    \end{flushleft}

\vfill
14 января 2025

\end{center}
\end{titlepage}

%%%%%%%%%%%%%%%%%%%
\renewcommand{\abstractname}{\Large Аннотация}
\begin{abstract}
В статье рассматриваются методы минимизации погрешностей при дискретной передаче данных между системами. Задача оптимального выбора интервала передачи формализована с использованием аксиоматического подхода, учитывающего ограничения на точность и сохранение данных. На основе введённой аксиоматики были разработаны два алгоритма, применимость которых подтверждена эмпирическим исследованием. Обсуждаются также возможности адаптации методов для специфических случаев, включая квантовые системы, и предлагаются направления для дальнейших исследований.
\end{abstract}

%%%%%%%%%%%%%%%%%%%
\tableofcontents
\pagebreak
\renewcommand{\thesection}{\arabic{section}}

%%%%%%%%%%%%%%%%%%%
\section{Введение}
Рассмотрим следующую модель: система, которую далее будем называть \textbf{Отправитель}, находится в некотором состоянии $\bm{\gamma}$, общее число которых конечно. Данные о своём состоянии она передаёт в другую систему, которую будем называть \textbf{Получатель}. Отправитель может изменять своё состояние в зависимости от времени $\bm{t}$. Главной целью является минимизация потерь данных при передаче, чтобы максимально сохранить информацию о зависимости состояния от времени.

Для наглядности представим зависимость состояния от времени в виде графика, где ось абсцисс обозначает время, а ось ординат — тип состояния:
\begin{figure}[ht]
    \centering
    \begin{tikzpicture}
        \begin{axis}[
            xlabel={Время $t$},
            ylabel={Состояние},
            xmin=0, xmax=100,
            ymin=0, ymax=5,
            ytick={0,1,2,3,4},
            ymajorgrids=true,
            grid style=dashed,
            axis lines = middle,
            enlargelimits = true,
            x post scale = 2,
            xmajorticks=false
            ]

        \addplot[no markers] coordinates {(0,1)(20,1)};
        \addplot[no markers] coordinates {(20,4)(30,4)};
        \addplot[no markers] coordinates {(30,3)(80,3)};
        \addplot[no markers] coordinates {(80,2)(120,2)};

        \addplot[dashed, no markers] coordinates {(20,1)(20,4)};
        \addplot[dashed, no markers] coordinates {(30,4)(30,3)};
        \addplot[dashed, no markers] coordinates {(80,3)(80,2)};

        \node[below left] at (axis cs:0,0) {$0$};
        \end{axis}
    \end{tikzpicture}
    \caption{график зависимости состояния от времени.}
    \label{fig:example}
\end{figure}

На данном графике каждому моменту времени $t$ соответствует только одно состояние системы. Это свойство обеспечивает однозначность зависимости.
\label{sec:def}

Тривиальное решение задачи заключается в использовании непрерывного сигнала, который передаёт данные без потерь. Однако в некоторых системах (например, в квантовых) такое решение невозможно. Дискретность таких систем делает задачу менее очевидной. В данной статье основное внимание уделено разработке методов минимизации потерь информации для дискретных случаев.

%%%%%%%%%%%%%%%%%%%
\section{Метод минимизации суммы погрешностей (ММСП)}
\subsection{Аксиоматика}
Для решения задачи в условиях дискретности важно учитывать следующие положения:

\quad I. Если интервал передачи данных слишком велик, полученные результаты будут иметь значительную погрешность, что сведёт ценность данных к нулю.

\quad II. Если интервал передачи данных слишком мал, точность данных будет высокой, однако это создаст чрезмерную нагрузку на устройство и сеть передачи информации.

Также будем считать, что
\begin{enumerate}[label=\asbuk*),ref=\asbuk*]
    \item Важно сохранить информации как о виде состояния, так и его длительности;
    \item Различные виды состояния могут иметь разные периоды отправки сигнала;
    \item Отправитель не может менять период отправки в течение времени.
    \label{sec:c}
\end{enumerate}
Случай одинакового периода отправки для различных состояний является частным к рассматриваемому, поэтому решение для такой проблемы может быть легко получено, используя далее описанные методы.

%%%%%%%%%%%%%%%%%%%
\subsection{Математическое введение}
Совокупность состояния и его длительности будем называть промежутком. Чтобы сохранить информацию о каждом промежутке, необходимо минимизировать искажения вида и длительности каждого состояния. Следовательно, промежутки должны рассматриваться независимо друг от друга.

Пусть интервал отправки, обозначаемый $\bm{\tau}$, зависит от длительности промежутков данного состояния. Мультимножество всех промежутков одного состояния обозначим как $\bm{T}$:

Логично предположить, что интервал отправки, который далее будем обозначать $\bm{\tau}$, зависит от длительности промежутков для данного состояния. Мультимножество (множество, где могут быть повторения элементов) всех промежутков одного состояния обозначим как $\bm{T}$:
$$T=\{t_i\ |\ t_i\in \mathbb{Q}\}$$
Здесь $t_i$ — длительность промежутков, которые являются рациональными числами, поскольку любые данные имеют конечную точность. Мультимножество $T$ формируется из собранной базы данных.

Каждый промежуток $t_i$ можно выразить следующим образом:
$$t_i=\bm{m_i} \tau+\bm{\Delta t_i}\ (1)$$
где $m_i$ — число отправленных интервалов длиной $\tau$, а $\Delta t_i$ — разница между реальной длительностью $t_i$ и её отправленной частью, которую назовём \textbf{отклонением}.

Число $m_i$ можно представить в виде
$$m_i=\left[\frac{t_i}{\tau}\right]\ (2)$$
где квадратные скобки обозначают округление к ближайшему целому. Это округление минимизирует суммарную абсолютную погрешность, поскольку обеспечивает симметрию: одни промежутки округляются вниз, а другие — вверх. Альтернативные подходы, например, всегда округлять вниз или вверх, дают более высокую абсолютную погрешность.

Промежуток имеет наглядное представление:
\begin{figure}[ht]
    \centering
    \begin{tikzpicture}
        \draw (0, 0) -- (5, 0);
        \draw (0, 0.25) -- (0, -0.25);
        \draw (2, 0.25) -- (2, -0.25);
        \draw (4, 0.25) -- (4, -0.25);
        \draw (5, 0.25) -- (5, -0.25);

        \draw (1, 0) node[anchor=south]{$\tau$};
        \draw (3, 0) node[anchor=south]{$\tau$};
        \draw (4.5, 0) node[anchor=south]{$\Delta t_i$};
        \draw (2.5, -0.25) node[anchor=north]{$t_i$};


    \end{tikzpicture}
    \caption{промежуток.}
\end{figure}

%%%%%%%%%%%%%%%%%%%
\subsection{Сильное условие нахождения интервала}\label{sec:GCF}

Условие для достижения искомого результата формулируется как: 

$$\Delta t_i=0,\ \forall t_i\ (3)$$

При соблюдении этого условия формула (1) принимает следующий вид:

$$m_i=\frac{t_i}{\tau}\ (4)$$

Объединяя уравнения (2) и (4), мы имеем:

$$\frac{t_i}{\tau}=\left[ \frac{t_i}{\tau} \right] (5)$$

Условие (5) должно выполняться для каждого из обозначенных интервалов, что указывает на то, что $\tau$ является наибольшим общим делителем (НОД) мультимножества $T$, так как он гарантирует неизменность значения $\frac{t_i}{\tau}$ при округлении. Хотя традиционно НОД для рациональных чисел не рассматривается, мы можем домножить каждое $t_i$ на одно общее число $\bm{a}$, устранив тем самым дробные части и позволяя работать с $at_i$ (при этом $\tau$ также домножится на это число):

$$\frac{at_i}{a\tau}=\left[ \frac{at_i}{a\tau} \right] (6)$$

Это возможно сделать, поскольку исходные значения имеют ограниченную точность. При этом важно, чтобы коэффициент $a$ был минимально возможным, чтобы не допустить возникновения нежелательных общих множителей, общих для всех чисел $at_i$ одновременно. В этом случае $\tau$ определяется следующим уравнением:

$$a\tau=\text{НОД}(T)$$
$$\tau=\frac{\text{НОД}(T)}{a}$$

Поясним, что для произвольного набора данных $T$, НОД мультимножества $aT=\{at_i\}$ стремится к 1 по мере увеличения числа и разнообразия элементов, что приводит в общем случае к следующему выражению для $\tau$:

$$\tau=\frac{1}{a}$$

Данный метод нахождения периода отправки определяет слишком малый интервал, что вызывает проблему II, особенно если принять во внимание, что значение $a$ превышает единицу.

%%%%%%%%%%%%%%%%%%%
\subsection{Слабое условие нахождения интервала}
Поскольку обеспечить условие (3) без противоречия условиям I и II не удалось, то потребуем минимальности отклонения:
$$|\Delta t_i|\rightarrow|\Delta t_i^{min}|\ (7)$$
Так как $\Delta t_i$ является функцией от $\tau$, то минимум можем найти используя производные:
$$|\Delta t_i|'_{\tau}=0\text{ — условие экстремума}\ (8)$$ 
$$|\Delta t_i|''_{\tau}>0\text{ — условие минимума}\ (9)$$

В равенстве (1) $\Delta t_i=t_i-\left[\frac{t_i}{\tau}\right]\tau$ фигурирует функция округления, поэтому для раскрытия условия (8) необходимо знать её производную. Для этого рассмотрим график функции округления $\left[\frac{t_i}{\tau}\right]=f(\frac{t_i}{\tau})$:

\begin{figure}[ht]
    \centering
    \begin{tikzpicture}
        \begin{axis}[
            xlabel={$\frac{t_i}{\tau}$},
            ylabel={$\left[\frac{t_i}{\tau}\right]$},
            xmin=0, xmax=5,
            ymin=0, ymax=5,
            xtick={0,1,2,3,4,5},
            ytick={0,1,2,3,4,5},
            ymajorgrids=true,
            grid style=dashed,
            axis lines = middle,
            enlargelimits = true,
        ]
        
        \addplot[no markers] coordinates {(0,0)(0.5,0)};
        \addplot[no markers] coordinates {(0.5,1)(1.5,1)};
        \addplot[no markers] coordinates {(1.5,2)(2.5,2)};
        \addplot[no markers] coordinates {(2.5,3)(3.5,3)};
        \addplot[no markers] coordinates {(3.5,4)(4.5,4)};
        \addplot[no markers] coordinates {(4.5,5)(5.5,5)};
        
        \addplot[dashed, no markers] coordinates {(0.5,0)(0.5,1)};
        \addplot[dashed, no markers] coordinates {(1.5,1)(1.5,2)};
        \addplot[dashed, no markers] coordinates {(2.5,2)(2.5,3)};
        \addplot[dashed, no markers] coordinates {(3.5,3)(3.5,4)};
        \addplot[dashed, no markers] coordinates {(4.5,4)(4.5,5)};
        
        \addplot[only marks, mark=*, mark options={scale=0.75, fill=white}] coordinates {
            (0.5,0)(1.5,1)(2.5,2)(3.5,3)(4.5,4)
        };
        \addplot[only marks, mark=*, mark options={scale=0.75}] coordinates {
            (0.5,1)(1.5,2)(2.5,3)(3.5,4)(4.5,5)
        };

        \node[below left] at (axis cs:0,0) {$0$};
        \end{axis}
    \end{tikzpicture}
    \caption{функция округления.}
\end{figure}

Но в аргументе функции $\tau$ стоит в знаменателе, поэтому для нахождения производной построим график от обратного аргумента $\left[\frac{t_i}{\tau}\right]=f(\frac{\tau}{t_i})$:
\newpage
\begin{figure}[ht]
    \centering
    \begin{tikzpicture}
        \begin{axis}[
            xlabel={$\frac{\tau}{t_i}$},
            ylabel={$\left[\frac{t_i}{\tau}\right]$},
            xmin=0, xmax=2,
            ymin=0, ymax=5,
            xtick={0,2/8,4/8,6/8,8/8,10/8,12/8,14/8,16/8},
            ytick={0,1,2,3,4,5},
            ymajorgrids=true,
            grid style=dashed,
            axis lines = middle,
            enlargelimits = true,
            xticklabels={$\frac{1}{8}$,$\frac{2}{8}$,$\frac{4}{8}$,$\frac{6}{8}$,$\frac{8}{8}$,$\frac{10}{8}$,$\frac{12}{8}$,$\frac{14}{8}$,$\frac{16}{8}$},
            y label style={at={(axis description cs:0.1,1)},anchor=south east}
        ]
        
        \addplot[no markers] coordinates {(3,0)(2/1,0)};
        \addplot[no markers] coordinates {(2/1,1)(2/3,1)};
        \addplot[no markers] coordinates {(2/3,2)(2/5,2)};
        \addplot[no markers] coordinates {(2/5,3)(2/7,3)};
        \addplot[no markers] coordinates {(2/7,4)(2/9,4)};
        \addplot[no markers] coordinates {(2/9,5)(2/11,5)};
        
        \addplot[dashed, no markers] coordinates {(2/1,0)(2/1,1)};
        \addplot[dashed, no markers] coordinates {(2/3,1)(2/3,2)};
        \addplot[dashed, no markers] coordinates {(2/5,2)(2/5,3)};
        \addplot[dashed, no markers] coordinates {(2/7,3)(2/7,4)};
        \addplot[dashed, no markers] coordinates {(2/9,4)(2/9,5)};
        \addplot[dashed, no markers] coordinates {(2/11,5)(2/11,6)};
        
        \addplot[only marks, mark=*, mark options={scale=0.75, fill=white}] coordinates {
            (2,0)(2/3,1)(2/5,2)(2/7,3)(2/9,4)(2/11,5)
        };
        \addplot[only marks, mark=*, mark options={scale=0.75}] coordinates {
            (2/1,1)(2/3,2)(2/5,3)(2/7,4)(2/9,5)
        };
        \node[below left] at (axis cs:0,0) {$0$};
        \end{axis}
    \end{tikzpicture}
    \caption{зависимость значения функции округления от $\frac{\tau}{t_i}$.}
\end{figure}

Как видно, график в любой точке $\frac{\tau}{t_i}$ параллелен оси абцисс, а значит производная от функции округления в любой точке на области определения равна 0. Теперь можно раскрыть формулу (8):
$$\left|t_i-\left[\frac{t_i}{\tau}\right]\tau\right|'_\tau=-\frac{t_i-\left[\frac{t_i}{\tau}\right]\tau}{|t_i-\left[\frac{t_i}{\tau}\right]\tau|}\left(\left[\frac{t_i}{\tau}\right]'_\tau\tau+\left[\frac{t_i}{\tau}\right]\tau'_\tau\right)=$$
$$=-\frac{t_i-\left[\frac{t_i}{\tau}\right]\tau}{|t_i-\left[\frac{t_i}{\tau}\right]\tau|}\left(0+1\cdot\left[\frac{t_i}{\tau}\right]\right)=-\frac{t_i-\left[\frac{t_i}{\tau}\right]\tau}{|t_i-\left[\frac{t_i}{\tau}\right]\tau|}\left[\frac{t_i}{\tau}\right]$$
Множитель $\frac{t_i-\left[\frac{t_i}{\tau}\right]\tau}{\left|t_i-\left[\frac{t_i}{\tau}\right]\tau\right|}$ не может равняться нулю, так как в единственной подозрительной на это точке он неопределён, то
$$\left[\frac{t_i}{\tau}\right]=0\ (10)$$
Данное условие можно переписать в виде
$$\frac{t_i}{\tau}<\frac{1}{2}\Leftrightarrow\tau>2t_i$$

что верно для всех промежутков только при $\tau>2\bm{\max(T)}$, где $\max(T)$ — максимальный элемент из $T$.

Данное значение $\tau$ создаёт проблему I, поэтому введём другое условие, основанное на минимизации суммы отклонений (7):
$$\bm{\Phi(\tau)}=\sum_{i=1}^n|\Delta t_i|\rightarrow\left(\sum_{i=1}^n|\Delta t_i|\right)^{min}\ (11)$$
Это условие целесообразно для решения поставленной задачи, поскольку все промежутки равноправны. Важно отметить, что условие (11) не является эквивалентным условию (7), так как (7) минимизирует каждое отдельное отклонение $|\Delta t_i|$, тогда как (11) минимизирует только их сумму. Это означает, что некоторые значения $|\Delta t_i|$ могут не достигать своих индивидуальных минимумов, если другие значения $|\Delta t_i|$ компенсируют это.

Применим тот же метод, что и для (7):
$$\Phi'_\tau(\tau)=\left(\sum_{i=1}^n|\Delta t_i|\right)'_\tau=\sum_{i=1}^n|\Delta t_i|'_\tau=-\sum_{i=1}^n\frac{t_i-\left[\frac{t_i}{\tau}\right]\tau}{|t_i-\left[\frac{t_i}{\tau}\right]\tau|}\left[\frac{t_i}{\tau}\right]$$
$$\sum_{i=1}^n\frac{t_i-\left[\frac{t_i}{\tau}\right]\tau}{|t_i-\left[\frac{t_i}{\tau}\right]\tau|}\left[\frac{t_i}{\tau}\right]=0\ (12)$$

Данное уравнение уже имеет нетривиальные решения без противоречия условиям I и II, поэтому проверим условие (9):
$$\Phi''_\tau(\tau)=\left(\sum_{i=1}^n|\Delta t_i|\right)''_\tau>0$$
$$\left(\sum_{i=1}^n|\Delta t_i|\right)''_\tau=\left(-\sum_{i=1}^n\frac{t_i-\left[\frac{t_i}{\tau}\right]\tau}{|t_i-\left[\frac{t_i}{\tau}\right]\tau|}\left[\frac{t_i}{\tau}\right]\right)'_\tau=$$
$$=-\sum_{i=1}^n\left(\frac{t_i-\left[\frac{t_i}{\tau}\right]\tau}{|t_i-\left[\frac{t_i}{\tau}\right]\tau|}\left[\frac{t_i}{\tau}\right]'_\tau+\left(\frac{t_i-\left[\frac{t_i}{\tau}\right]\tau}{|t_i-\left[\frac{t_i}{\tau}\right]\tau|}\right)'_\tau\left[\frac{t_i}{\tau}\right]\right)=$$
$$=-\sum_{i=1}^n\left(0+\frac{(t_i-\left[\frac{t_i}{\tau}\right]\tau)'_\tau\cdot|t_i-\left[\frac{t_i}{\tau}\right]\tau|-(t_i-\left[\frac{t_i}{\tau}\right]\tau)|t_i-\left[\frac{t_i}{\tau}\right]\tau|'_\tau}{|t_i-\left[\frac{t_i}{\tau}\right]\tau|^2}\left[\frac{t_i}{\tau}\right]\right)=$$
$$=-\sum_{i=1}^n\frac{(-\left[\frac{t_i}{\tau}\right])\cdot|t_i-\left[\frac{t_i}{\tau}\right]\tau|+(t_i-\left[\frac{t_i}{\tau}\right]\tau)\frac{t_i-\left[\frac{t_i}{\tau}\right]\tau}{|t_i-\left[\frac{t_i}{\tau}\right]\tau|}\left[\frac{t_i}{\tau}\right]}{|t_i-\left[\frac{t_i}{\tau}\right]\tau|^2}\left[\frac{t_i}{\tau}\right]=$$
$$=\sum_{i=1}^n\frac{\left[\frac{t_i}{\tau}\right]|t_i-\left[\frac{t_i}{\tau}\right]\tau|-\frac{(t_i-\left[\frac{t_i}{\tau}\right]\tau)^2}{|t_i-\left[\frac{t_i}{\tau}\right]\tau|}\left[\frac{t_i}{\tau}\right]}{|t_i-\left[\frac{t_i}{\tau}\right]\tau|^2}\left[\frac{t_i}{\tau}\right]=\left\{(x)^2=|x|^2\right\}$$
$$=\sum_{i=1}^n\frac{|t_i-\left[\frac{t_i}{\tau}\right]\tau|-\frac{|t_i-\left[\frac{t_i}{\tau}\right]\tau|^2}{|t_i-\left[\frac{t_i}{\tau}\right]\tau|}}{|t_i-\left[\frac{t_i}{\tau}\right]\tau|^2}\left[\frac{t_i}{\tau}\right]^2=$$
$$=\sum_{i=1}^n\frac{|t_i-\left[\frac{t_i}{\tau}\right]\tau|-|t_i-\left[\frac{t_i}{\tau}\right]\tau|}{|t_i-\left[\frac{t_i}{\tau}\right]\tau|^2}\left[\frac{t_i}{\tau}\right]^2=0$$
$$\Phi''_\tau(\tau)=0\ (13)$$

На последнем этапе исчезло деление модулей, которое создаёт неопределённость в некоторых точках. Вне этих точек на всей области определения вторая производная равна нулю, что означает, что все точки являются только точками перегиба. Таким образом, найти минимумы через производные не получится.

На нижеприведённом рисунке изображён график функции $\Phi(\tau)$ для \hyperref[sec:wait]{режима ожидания}:
\begin{figure}[ht]
    \centering
    \begin{tikzpicture}
        \begin{axis}[
            xlabel={$\frac{\tau}{t_i}$},
            ylabel={$\left[\frac{t_i}{\tau}\right]$},
            xmin=0, xmax=4,
            ymin=0, ymax=0.3,
            xtick={0.5, 1.0, 1.5, 2.0, 2.5, 3.0, 3.5, 4.0},
            ytick={0,0.1,0.2,0.3},
            axis lines = middle,
            enlargelimits = true,
            x post scale=2,
            xticklabel style={/pgf/number format/.cd, fixed, fixed zerofill, precision=1}
            ]
        
        \addplot[ultra thin, no markers] file {phi.txt};

        \node[below left] at (axis cs:0,0) {$0$};
        \end{axis}
    \end{tikzpicture}
    \caption{график функции $\Phi(\tau)$.}
\end{figure}

%%%%%%%%%%%%%%%%%%%
\subsection{\textbf{Алгоритм}}

Поскольку производные не помогли, найдём минимумы с помощью программирования: будем вычислять значения $\Phi(\tau)$ с определённым шагом для $\tau$ на заданном промежутке и искать минимальное значение. Затем уточним это значение, уменьшая шаг, пока не достигнем желаемой точности.

Промежуток поиска минимального значения определим следующим образом: максимальное значение $\tau$ не может превышать $\min(T)$, чтобы сохранить информацию о всех интервалах. Нижняя граница определяется максимальной по модулю абсолютной погрешностью данных в $T$ ($\Delta_{\max} t$), так как меньшие значения $\tau$ будут неверными с точки зрения погрешности. Окончательное значение $\tau$ округляется с учётом абсолютной погрешности исходных данных.

$$\Delta_{\max} t \leq \tau \leq \min(T)\ (14)$$

Таким образом, получаем следующий алгоритм:
\begin{enumerate}
    \item Получить выборку $T$ из промежутков одинакового типа;
    \item Вычислить сумму $\sum_{i=1}^n \left|t_i - \left[\frac{t_i}{\tau}\right]\tau\right|$ ($n$ — число элементов выборки) с определённым шагом на промежутке (14);
    \item Найти минимум суммы на данном интервале и соответствующее значение $\tau$;
    \item Уточнить значение $\tau$ в окрестности минимума при необходимости;
    \item Значение $\tau$, соответствующее минимальному значению суммы, является искомым.
\end{enumerate}

%%%%%%%%%%%%%%%%%%%
\section{Центрированный наибольший общий делитель (ЦНОД)}
\subsection{Идея}

Вернёмся к \hyperref[sec:GCF]{сильному условию нахождения интервала}. В процессе размышлений о наибольшем общем делителе (НОД) был найден альтернативный алгоритм, позволяющий определить $\tau$ без противоречия условиям I и II. Обозначим $a t_i$ как $\bm{d_i}$. Тогда произвольный промежуток $d$ можно представить через его простые множители:

$$d = \bm{p_1^{k_1}} \cdot ... \cdot \bm{p_r^{k_r}}\ (15),$$

где $p$ — простой множитель, $k$ — кратность множителя, а $r$ — количество уникальных множителей.

Посчитаем количество всех простых делителей во всех промежутках без учёта кратности, то есть:

$$d_i: p_{i1}, ..., p_{i r_i}$$
$$\bm{P} = \{2: \bm{l_2}, 3: \bm{l_3}, ...,p_{\bm{N}}: \bm{l_N}\},$$

где $P$ — множество простых множителей из выборки, $N$ — их количество, а $l$ — число вхождений определённого множителя в различных промежутках.

Кратность не учитывается, так как при отправке сигнала промежуток делится на $\tau$ только один раз.

В итоге идеальный интервал определяется по следующей формуле:

$$\tau = \frac{\sum_{i=2}^N l_i \cdot p_i}{\sum_{j=2}^N l_j}\ (16)$$

Эта формула полностью аналогична координате центра масс системы материальных точек, что и дало название алгоритму:

$$x_C = \frac{\sum_{i=1}^n m_i \cdot x_i}{\sum_{i=1}^n m_i}$$

\subsection{\textbf{Алгоритм}}

Таким образом, для определения идеального интервала получается следующий алгоритм:
\begin{enumerate}
    \item Получить выборку $T$ из промежутков одинакового типа;
    \item Сформировать выборку $aT$;
    \item Найти простые делители всех промежутков $d$;
    \item Посчитать количество раз, с которым каждый простой множитель встречается в промежутках;
    \item Рассчитать $\tau$ по формуле (16).
\end{enumerate}

%%%%%%%%%%%%%%%%%%%
\section{Эмпирическое сравнение дискретных алгоритмов}

Для проверки теоретических результатов в качестве частного случая была рассмотрена задача передачи данных с телефона на другое устройство о деятельности пользователя в течение одного цикла разрядки (см. \hyperref[sec:attach]{Приложение}). Данные собирались на одном устройстве в течение 27 дней. Типичный цикл представлен на рисунке:
\label{sec:wait}

\begin{center}
\begin{BVerbatim}
33м 37в 88о 1з 14м 3з 94о 75м 14в 2з 476о 13м
\end{BVerbatim}
\end{center}

Здесь числа обозначают минуты беспрерывного использования смартфона в одном режиме. Учитывались четыре режима: «о» — ожидание, «з» — звонок, «м» — прослушивание музыки, «в» — просмотр видео. Для каждого режима вычислялся свой идеальный интервал по данным всех промежутков одного вида из всех циклов.

Введём погрешность передачи отдельных промежутков, которую будем называть дифференциальной:

$$\Delta t_{\text{д}} = |\Delta t_i|,\ \eps_{\text{д}} = \frac{|\Delta t_i|}{t_i}$$

Для погрешности передачи целых циклов, обозначаемой как интегральная, используем следующие формулы:

$$\Delta t_{\text{и}} = \sum_{i=1}^{\bm{q}} |\Delta t_i|,\ \eps_{\text{и}} = \frac{\sum_{i=1}^{\bm{q}} |\Delta t_i|}{\sum_{i=1}^{\bm{q}} t_i}$$

где $q$ — число промежутков в данном цикле.

Параметры для алгоритмов:
\begin{itemize}
    \item ЦНОД: $a = 1$;
    \item ММСП: $\Delta_{\max} t$ оценивается как $1$.
\end{itemize}


Результаты работы алгоритмов для вышеупомянутой выборки: 

\begin{minipage}[b]{0.45\textwidth}
    \centering ЦНОД

    Дифференциальная:
        $$\eps_{\min}=0\%, \Delta t_{\min}=0\ \text{мин}$$
        $$\eps_{\max}=1\cdot 10^2\%, \Delta t_{\max}=16\ \text{мин}$$
        $$\eps_{\text{ср.}}=11\%, \Delta t_{\text{ср.}}=5\ \text{мин}$$
    Интегральная:
        $$\eps_{\min}=1,6\%, \Delta t_{\min}=15\ \text{мин}$$
        $$\eps_{\max}=7\%, \Delta t_{\max}=68\ \text{мин}$$
        $$\eps_{\text{ср.}}=4\%, \Delta t_{\text{ср.}}=37\ \text{мин}$$
\end{minipage}
\hfill
\begin{minipage}[b]{0.45\textwidth}
    \centering ММСП

    Дифференциальная:
    $$\eps_{\min}=0\%, \Delta t_{\min}=0\ \text{мин}$$
    $$\eps_{\max}=2\cdot10^{1}\%, \Delta t_{\max}=2\ \text{мин}$$
    $$\eps_{\text{ср.}}=1,5\%, \Delta t_{\text{ср.}}=1\ \text{мин}$$
    Интегральная:
    $$\eps_{\min}=0,07\%, \Delta t_{\min}=1\ \text{мин}$$
    $$\eps_{\max}=0,8\%, \Delta t_{\max}=8\ \text{мин}$$
    $$\eps_{\text{ср.}}=0,5\%, \Delta t_{\text{ср.}}=4\ \text{мин}$$
\end{minipage}

\vspace{0.5cm}
Интервалы отправки для данной выборки:
\vspace{0.5cm}

\begin{minipage}[b]{0.45\textwidth}
    \centering ЦНОД:
    \begin{itemize}
        \item Ожидание — 33 (32,8) мин
        \item Видео — 11 (10,9) мин
        \item Музыка — 13 (12,6) мин
        \item Звонок — 3 (2,6) мин
    \end{itemize}\end{minipage}
\hfill
\begin{minipage}[b]{0.45\textwidth}
    \centering ММСП:
    \begin{itemize}
        \item Ожидание — 1 (1,2) мин
        \item Видео — 5 (4,8) мин
        \item Музыка — 1 (1,4) мин
        \item Звонок — 1 (1,0) мин
    \end{itemize}
\end{minipage}

Значения погрешностей рассчитаны с одной запасной цифрой для $\tau$ (значение с запасной цифрой указано в скобках). Погрешности исходных данных при расчёте не учитывались.

Для оценки качества работы алгоритмов удобно ввести величину, которую назовём \textbf{добротностью}:

$$\zeta = \frac{|\Delta t|}{\tau} \cdot 100\% \ (17)$$

где $\Delta t$ — отклонение для одного типа данных.

Чем меньше значение добротности, тем лучше алгоритм, поскольку интервал отправки данных становится настолько большим, что отклонение становится незначительным, эффективно перекрывая его и минимизируя временные "провалы" 
 при отправке. Добротность всегда является дифференциальной, а также может классифицироваться как минимальная, максимальная и средняя. Найдём средние добротности для всех типов деятельности для данной выборки:
\vspace{0.5cm}

\begin{minipage}[b]{0.45\textwidth}
    \centering ЦНОД:
    \begin{itemize}
        \item Ожидание — 26,7 \%
        \item Видео — 23,5 \%
        \item Музыка — 24,4 \%
        \item Звонок — 27,7 \%
    \end{itemize}\end{minipage}
\hfill
\begin{minipage}[b]{0.45\textwidth}
    \centering ММСП:
    \begin{itemize}
        \item Ожидание — 25,1 \%
        \item Видео — 23,5 \%
        \item Музыка — 25,8 \%
        \item Звонок — 0,0 \%
    \end{itemize}
\end{minipage}

Как видно из вышеприведённых результатов, добротность алгоритмов практически одинакова, за исключением режима звонка, поскольку минимальное значение совпало с абсолютной погрешностью. Учитывая тот факт, что интервалы отправки данных у ЦНОД заметно больше, можно сделать вывод, что данный алгоритм является более совершенным, хотя и не был получен строгим выводом уравнений. Однако такое сравнение не совсем корректно — при использовании ЦНОД некоторые промежутки могут пропасть, в то время как верхняя граница ММСП ограничена минимальным промежутком, что не позволяет найти достаточно большой интервал отправки.

%%%%%%%%%%%%%%%%%%%%%
\section{Эмпирическое сравнение дискретных алгоритмов с модификацией ММСП}
Для более корректного сравнения алгоритмов сместим верхнюю границу ММСП до среднего арифметического для данной выборки:

$$\Delta_{\max} t \le \tau \le t_{\text{ср.}}$$

Все остальные параметры оставим теми же.

Результаты работы алгоритмов (значения для ЦНОД дублированы для наглядного сравнения):

\begin{minipage}[b]{0.45\textwidth}
    \centering ЦНОД

    Дифференциальная:
        $$\eps_{\min}=0\%, \Delta t_{\min}=0\ \text{мин}$$
        $$\eps_{\max}=1\cdot 10^2\%, \Delta t_{\max}=16\ \text{мин}$$
        $$\eps_{\text{ср.}}=11\%, \Delta t_{\text{ср.}}=5\ \text{мин}$$
    Интегральная:
        $$\eps_{\min}=1,6\%, \Delta t_{\min}=15\ \text{мин}$$
        $$\eps_{\max}=7\%, \Delta t_{\max}=68\ \text{мин}$$
        $$\eps_{\text{ср.}}=4\%, \Delta t_{\text{ср.}}=37\ \text{мин}$$
\end{minipage}
\hfill
\begin{minipage}[b]{0.45\textwidth}
    \centering ММСП

    Дифференциальная:
    $$\eps_{\min}=0\%, \Delta t_{\min}=0\ \text{мин}$$
    $$\eps_{\max}=1\cdot 10^2\%, \Delta t_{\max}=17\ \text{мин}$$
    $$\eps_{\text{ср.}}=6\%, \Delta t_{\text{ср.}}=4\ \text{мин}$$
    Интегральная:
    $$\eps_{\min}=1,2\%, \Delta t_{\min}=12\ \text{мин}$$
    $$\eps_{\max}=6\%, \Delta t_{\max}=47\ \text{мин}$$
    $$\eps_{\text{ср.}}=3\%, \Delta t_{\text{ср.}}=28 \text{мин}$$
\end{minipage}

\vspace{0.5cm}
Интервалы отправки для данной выборки:
\vspace{0.5cm}

\begin{minipage}[b]{0.45\textwidth}
    \centering ЦНОД:
    \begin{itemize}
        \item Ожидание — 33 (32,8) мин
        \item Видео — 11 (10,9) мин
        \item Музыка — 13 (12,6) мин
        \item Звонок — 3 (2,6) мин
    \end{itemize}\end{minipage}
\hfill
\begin{minipage}[b]{0.45\textwidth}
    \centering ММСП:
    \begin{itemize}
        \item Ожидание — 35 (35,0) мин
        \item Видео — 5 (4,8) мин
        \item Музыка — 1 (1,4) мин
        \item Звонок — 1 (1,3) мин
    \end{itemize}
\end{minipage}

Добротности:
\vspace{0.5cm}

\begin{minipage}[b]{0.45\textwidth}
    \centering ЦНОД:
    \begin{itemize}
        \item Ожидание — 26,7 \%
        \item Видео — 23,5 \%
        \item Музыка — 24,4 \%
        \item Звонок — 27,7 \%
    \end{itemize}\end{minipage}
\hfill
\begin{minipage}[b]{0.45\textwidth}
    \centering ММСП:
    \begin{itemize}
        \item Ожидание — 25,6 \%
        \item Видео — 23,5 \%
        \item Музыка — 25,8 \%
        \item Звонок — 22,3 \%
    \end{itemize}
\end{minipage}

Как видно из вышеприведённых результатов, добротность алгоритмов снова практически совпала. Средние дифференциальные и интегральные погрешности у ММСП оказались немного ниже, а добротность алгоритмов колеблется около общего среднего значения. Из этого следует вывод, что нет однозначного перевеса точности в сторону какого-либо алгоритма, и для выбора подходящего алгоритма в конкретном случае необходимо проводить эмпирические исследования, пока не будет установлено теоретическое обоснование разницы в точности.

%%%%%%%%%%%%%%%%%%%
\section{Итоги}
\subsection{Вывод}
В данной статье рассмотрены методы минимизации погрешностей при передаче данных с использованием алгоритмов ММСП и ЦНОД. На основе предложенной аксиоматики и математических моделей были разработаны алгоритмы, которые обеспечивают оптимальные интервалы передачи данных. Эмпирическое сравнение показало, что оба алгоритма демонстрируют схожие результаты в плане точности, с небольшими различиями в некоторых режимах. 

Модификация ММСП позволила добиться более гибких результатов, что указывает на необходимость дальнейших исследований для адаптации методов к специфическим требованиям. В целом, предложенные алгоритмы показали свою применимость в задачах минимизации потерь данных, и их эффективность подтверждена как теоретически, так и эмпирически.

%%%%%%%%%%%%%%%%%%%
\subsection{Развитие идеи}
В данной главе рассматриваются дополнительные аспекты проблемы, включая возможные изменения в аксиоматике, альтернативные функции округления и особенности квантовых систем, которые могут повлиять на выбор методов для нахождения оптимальных интервалов отправки данных. Эти идеи открывают новые направления для дальнейших исследований.

\subsubsection{Вариация аксиоматики}
Рассмотрим случай, когда Отправитель может изменять интервал во времени, то есть не выполняется \hyperref[sec:c]{пункт из аксиоматики}. Это означает, что пакеты данных будут поступать к Получателю с разной периодичностью. Очевидно, что наилучшим вариантом будет отправлять данные сразу же после переключения состояния, так как таким образом интервал будет максимальным и без потери информации.

\subsubsection{Другая функция округления}
Ещё раз вернёмся к \hyperref[sec:GCF]{сильному условию нахождения интервала}. Вместо округления к ближайшему целому можно построить собственную функцию на основе данной выборки $T$. Эту функцию обозначим как $\langle\frac{t_i}{\tau}\rangle$. Для создания этой функции нанесём точки, полученные с помощью округления к ближайшему целому:

\begin{figure}[ht]
    \centering
    \begin{tikzpicture}
        \begin{axis}[
            xlabel={$\frac{t_i}{\tau}$},
            ylabel={$\left[\frac{t_i}{\tau}\right]$},
            ymajorgrids=true,
            grid style=dashed,
            axis lines = middle,
            enlargelimits = true,
            x post scale=2,
            xticklabel style={/pgf/number format/.cd, fixed, fixed zerofill, precision=1}
            ]

            \addplot[ultra thin, only marks, mark=diamond*, mark options={scale=0.75}] file {round.txt};

        \node[below left] at (axis cs:0,0) {$0$};
        \end{axis}
    \end{tikzpicture}
    \caption{график округления к ближайшему целому.}
\end{figure}
\newpage
где $T = {4, 17, 28, 43, 44, 90, 94, 94, 117, 122}$ с $\tau = 32,8$ (часть режима ожидания и алгоритм ЦНОД).

Теперь проведём вертикали посередине между точками, между которыми величина абсциссы изменяется, и построим линии, параллельные оси абсцисс, от точек в обе стороны до вертикалей. Если между вертикальными составляющими каких-либо точек расстояние больше 1, то необходимо предварительно построить промежуточные точки, которые равномерно заполнят промежуточное пространство (обозначены крестиком):

\begin{figure}[ht]
    \centering
    \begin{tikzpicture}
        \begin{axis}[
            xlabel={$\frac{t_i}{\tau}$},
            ylabel={$\langle\frac{t_i}{\tau}\rangle$},
            xmin=0, xmax=4,
            ymin=0, ymax=4,
            ymajorgrids=true,
            grid style=dashed,
            axis lines = middle,
            enlargelimits = true,
            x post scale=2,
            xticklabel style={/pgf/number format/.cd, fixed, fixed zerofill, precision=1}
            ]

        \addplot[ultra thin, only marks, mark=diamond*, mark options={scale=0.75}] file {round.txt};
        
        \addplot[dashed, no markers] coordinates {(0.32,0)(0.32,5)};
        \addplot[dashed, no markers] coordinates {(1.692,0)(1.692,5)};
        \addplot[dashed, no markers] coordinates {(2.393,0)(2.393,5)};
        \addplot[dashed, no markers] coordinates {(3.217,0)(3.217,5)};
        \addplot[only marks, mark=x, mark options={scale=1.5}] coordinates {(2.043,2)};

        \addplot[no markers] coordinates {(0,0)(0.32,0)};
        \addplot[no markers] coordinates {(0.32,1)(1.692,1)};
        \addplot[no markers] coordinates {(1.692,2)(2.393,2)};
        \addplot[no markers] coordinates {(2.393,3)(3.217,3)};
        \addplot[no markers] coordinates {(3.217,4)(5,4)};

        \addplot[only marks, mark=*, mark options={scale=0.75, fill=white}] coordinates {
            (0.32,0)(1.692,1)(2.393,2)(3.217,3)
        };
        \addplot[only marks, mark=*, mark options={scale=0.75}] coordinates {
            (0.32,1)(1.692,2)(2.393,3)(3.217,4)
        };

        \node[below left] at (axis cs:0,0) {$0$};
        \end{axis}
    \end{tikzpicture}
    \caption{график новой функции округления.}
\end{figure}

Таким образом, получена новая функция округления, которая может уменьшить погрешность этой операции, так как результат будет концентрироваться ближе к точкам исходных данных. Хотя физический смысл такой операции не ясен, эта функция округления может найти своё применение.

\subsubsection{Квантовые системы и неоднозначность}
Как известно, квантовые системы характеризуются дискретностью энергии. Это обстоятельство требует отхода от непрерывных алгоритмов вследствие их неприменимости и обращения к дискретным методам, которые во многом были описаны в этой работе. Однако для создания прикладных способов передачи информации в квантовых системах необходимо провести отдельное исследование.

Также стоит отметить, что для описания изменения квантового состояния от времени может понадобиться отказаться от \hyperref[sec:def]{однозначности}: одной временной точке может соответствовать множество состояний $\gamma$ Отправителя, вплоть до бесконечности. Это положение требует разработки отдельных методов для решения задачи.

%%%%%%%%%%%%%%%%%%%
\section{Приложение}
\label{sec:attach}
\begin{verbatim}
    64в 383о 64в 101м 52в 236о
    35м 390о 124в 15м 205о 17в 19м 121в 11м
    30в 15м 495о 50в 69м 94о 23м 196о
    18м 28в 345о 197м 406о
    33м 37в 88о 1з 14м 3з 94о 75м 14в 2з 476о 13м
    715о 57в 72о 14м 35в 73о
    16м 20в 371о 36в 89м 227о 14в 60о 82м 4з 108м
    23в 20м 394о 76в 68м 192о 15м 20в
    38в 380о 28в 52о 1з 253о 11м 142о 72м 4о
    22м 23м 501о 66в 342о
    63м 233о 70в 282о 51в 226о
    59в 55м 104о 53в 47м 184о 42в 277о 5з
    34в 768о 1з 15м 106в 43о
    18м 928о
    7м 8з 430о 47в 50м 204о 14м 97в 123о
    65в 400о 120в 172о 106в 90о
    29в 34м 373о 53м 132в 117о 91в 150м
    26м 9в 383о 31в 4з 58в 445о 16м
    44в 446о 7м 69в 117о 79м 215о
    68в 236о 72в 26м 199о 35в 83о 66в 140о 30м
    17м 12в 496о 73м 244о 27в 30м 45о 20м
    81в 326о 25м 142в 44о 129м 33в 122о
    16м 45в 511о 72в 23м 246о 37в 17о 82м
    33в 36м 504о 21м 30в 372о
    18м 30в 561о 34в 25м 171о 61в 28о 52м
    40м 12в 484о 27в 77м 205о 88в 87о
    30м 35в 97о 16з 100о 10м 522о
\end{verbatim}

\end{document}
